\chapter{Appendix \thechapter}\label{appendixA}
%\chapter{Appendix A}\label{appendixA}
\appendixchaptername{Rubric Design and Prompt Templates}
      
\section{Rubric Development Process}

The rubrics used in this study were designed to ensure 
consistency and transparency in evaluating Thai written exam 
responses across Bloom’s taxonomy levels. The development 
followed a three-stage process:

\begin{enumerate}
    \item \textbf{Draft Generation:} Large Language Models (LLMs) such as ChatGPT were used to generate preliminary rubrics for each question based on Bloom’s cognitive level, target learning outcomes, and example student responses.
    \item \textbf{Expert Refinement:} Two professional Thai teachers reviewed and refined the generated rubrics. Ambiguous wording was clarified, and scoring descriptors were aligned with Thai pedagogical standards and curriculum guidelines.
    \item \textbf{Validation:} The finalized rubrics were tested on sample responses, and discrepancies between teacher and model evaluations were discussed and corrected before use in the main experiment.
\end{enumerate}

This iterative design ensured that rubrics provided both 
structured guidance for LLM scoring and reliable reference 
criteria for human teachers.

%%%%%%%%%%%%%%%%%%%%%%%%%%%%%%%%%%%%%%%%%%%%%%%%%%%%%%%%%%%%

\section{Example Rubric Table}

Each rubric defines four score levels (0–3) with clear 
performance descriptors and expected conceptual indicators. 
\rtab{tab:example_rubric} illustrates a 
representative rubric for the \textit{Understanding} level.

{
\centering
\begin{longtable}{c p{7cm} p{6cm}}
    \caption{Example Rubric for \textit{Understanding}-Level Question}
    \label{tab:example_rubric}
    \\

    % For Header
    \hline
    \textbf{Score} & \textbf{Description of Criteria} & \textbf{Example Keywords / Concepts} \\
    \hline
    \endfirsthead

    % For Overflow Header
    \hline
    \textbf{Score} & \textbf{Description of Criteria} & \textbf{Example Keywords / Concepts} \\
    \hline
    \endhead

    % For Footer
    \hline
    \endfoot

    3 & Response demonstrates complete comprehension, accurate explanation of concepts, and appropriate use of terminology in context. & Fully explains key ideas, correctly connects concepts, accurate interpretation \\
    2 & Response shows partial understanding; includes mostly correct ideas but lacks depth, clarity, or completeness. & Generally correct but missing details, partial explanation, limited linkage between concepts \\
    1 & Response reveals limited understanding or contains misconceptions about the main idea. & Vague or unclear explanation, minor conceptual errors, incomplete reasoning \\
    0 & Response is irrelevant, off-topic, or entirely incorrect. & Does not address the question, lacks essential content, incorrect reasoning \\
    
\end{longtable}
}

These rubrics were stored in Excel format for contextual 
retrieval during AI evaluation, allowing the LLMs to 
reference expected criteria dynamically when grading.

%%%%%%%%%%%%%%%%%%%%%%%%%%%%%%%%%%%%%%%%%%%%%%%%%%%%%%%%%%%%

\section{LLM Grading Prompt Templates}

Two main prompt templates were employed during the grading 
process: one for \textit{rubric-based evaluation} and another 
for \textit{non-rubric evaluation}. Both prompts were 
carefully structured to ensure consistency and minimize 
variance in language interpretation by the models.

\subsection*{Rubric-Based Prompt Template}

\begin{verbatim}
You are an expert Thai teacher.
Evaluate the following student's answer using the rubric below.
Assign a score (0-3) and explain your reasoning clearly.

[Question]: <question text>
[Student Answer]: <student answer>
[Rubric Criteria]:
Score 3: <criteria>
Score 2: <criteria>
Score 1: <criteria>
Score 0: <criteria>

Your Output Format:
Score: <number>
Reasoning: <concise explanation>
\end{verbatim}

\subsection*{Non-Rubric Prompt Template}

\begin{verbatim}
You are an expert Thai teacher.
Evaluate the student's written answer for correctness and completeness.
Assign a score between 0 and 3 and explain your reasoning.

[Question]: <question text>
[Student Answer]: <student answer>

Your Output Format:
Score: <number>
Reasoning: <concise explanation>
\end{verbatim}

%%%%%%%%%%%%%%%%%%%%%%%%%%%%%%%%%%%%%%%%%%%%%%%%%%%%%%%%%%%%

\section{Contextual Data Retrieval for Grading}

During rubric-based grading, contextual data were retrieved 
from the pre-validated Excel dataset. Each 
question entry included:

\begin{enumerate}
    \item Representative answers for each score level (0-3)
    \item Bloom’s taxonomy classification
    \item Rubric text and expected keywords
\end{enumerate}

This information was appended to the LLM input to provide 
contextual grounding and reduce ambiguity in interpretation.  
An example of contextual insertion is shown below.

\begin{verbatim}
[Contextual Reference Answers]
Score 3 Example: <full-mark answer>
Score 2 Example: <partial answer>
Score 1 Example: <minimal answer>
Score 0 Example: <irrelevant or incorrect answer>
\end{verbatim}

The combination of structured rubrics and contextual data 
helped standardize AI evaluation, reduce bias, and improve 
alignment with teacher grading in both rubric and non-rubric 
settings.
